\section{Related Work}
\label{sec:related}

\subsection{Verifiable rewards for code, with and without performance}

Post-training code models with reinforcement learning from verifiable
rewards - where an execution check rather than a human supplies the
signal~\cite{tulu3,deepseekr1}, optimized with GRPO~\cite{grpo} or its
successors~\cite{dapo,drgrpo,gspo} - is now standard. Several efforts include
a performance term. ACECode~\cite{acecode} fine-tunes with a dual-objective
rewarder that adds an efficiency component to a correctness base for Python
runtime. Kevin~\cite{kevin} rewards correctness plus a speedup ratio for GPU
kernels and documents the model lazily copying the reference implementation
rather than writing one. CUDA-L1~\cite{cudal1} rewards speedup and reports
timing exploits in which extra asynchronous streams mislead the timer.
CUDAPerf~\cite{cudaperf}, CUDA Agent~\cite{cudaagent}, and
MaxCode~\cite{maxcode} similarly fuse correctness with measured speedup.

Two lines are closest to ours and both come from the HPC community.
RLPF~\cite{rlpf} aligns code models to performance using a learned reward model
trained on fast/slow code pairs, with reported OpenMP gains of
1.9--4.5$\times$; the reward is predicted rather than measured, which is a
different design point with different failure modes, but the objective - get
the model to write faster parallel code - is the same one we are arguing for.
Real-machine reward GRPO~\cite{realmachine} post-trains an HPC code model using
measured GFLOPS on a supercomputer behind a staged penalty ladder, and is the
nearest existing instance of an execution-measured HPC reward. We differ in
object rather than in ambition: that work optimizes throughput on a kernel; we
ask why the correctness term cannot stand alone, and what structure follows.

Our contribution relative to all of the above is not the combination of
correctness and performance. It is (i) the argument that the correctness term's
optimum contains the serial program, which is what motivates a multiplicative
gate rather than a weighted sum, and (ii) an adversarial suite that turns
``which reward is better'' into a measurable question.

\subsection{Reward hacking and verifier gaming}

That verifiable rewards can be gamed is well documented. Recent work
distinguishes extensional from intensional verification and shows that
verifiers checking only extensional correctness admit false
positives~\cite{gamingverifiers}; others fuzz RLVR verifiers
directly~\cite{fuzzverifiers} and analyze reward hacking in rubric-based
RL~\cite{rubrichack}. Our serial projection is an extensional-verification
false positive of a particularly clean kind: it is not an exploit found by
search but a transform guaranteed to succeed on every correct program. \PG\
follows this literature's methodology - attack the reward, report what
survives - and instantiates it for parallel code.

\subsection{LLMs for parallel and HPC code}

\ParEval~\cite{pareval} is the standard benchmark, evaluating pass@$k$,
$\mathrm{speedup}_n@k$, and $\mathrm{efficiency}_n@k$ for model-generated
OpenMP, MPI, CUDA, and Kokkos, and documenting a large gap between models'
serial and parallel competence. Our capped efficiency
(Equation~\eqref{eq:efficiency}) is that benchmark's efficiency metric with a
unit cap, and we adopt it deliberately rather than inventing a variant.
\ParEval-Repo~\cite{parevalrepo} extends the setting to repository-level
translation. Model lines specialized for HPC include
HPC-Coder~\cite{hpccoder}, HPC-Coder-V2~\cite{hpccoderv2}, and
OMPGPT~\cite{ompgpt}, and PCEBench~\cite{pcebench} evaluates parallel code on
both correctness and performance. Work on whether models can predict parallel
code performance at all~\cite{predictperf} reinforces the premise that
performance is not implied by correctness.

Agentic HPC systems have proliferated: LASSI~\cite{lassi} builds a
compile-and-run self-correction loop for parallel-code translation,
MARCO~\cite{marco} and VibeCodeHPC~\cite{vibecodehpc} apply multi-agent and
auto-tuning approaches, UniPar~\cite{unipar} targets parallelization, and
ParEVO~\cite{parevo} guides an evolutionary loop with correctness
verification, race detection, and performance profiling - the closest existing
system to a two-signal loop, though the signals guide search rather than serve
as a reward.

\subsection{Functional-correctness and efficiency evaluation}

Code-generation evaluation has centered on functional correctness of small
self-contained problems~\cite{humaneval} and, at repository scale, on
SWE-bench~\cite{swebench}. EvalPlus~\cite{evalplus} showed that code passing a
visible test suite often fails under additional tests; our result is a
different failure of the same shape, along an axis no additional test can
reach. Efficiency-aware benchmarks - EffiBench~\cite{effibench},
ENAMEL~\cite{enamel} with its right-tail $\mathrm{eff}@k$ estimator, and
KernelBench~\cite{kernelbench} for GPU kernels - establish that measured
performance belongs in evaluation. We argue it belongs in the reward.

\subsection{Dynamic race detection}

A mature line of HPC tooling detects data races dynamically:
ThreadSanitizer~\cite{tsan} instruments memory accesses at scale, and
Archer~\cite{archer}, ROMP~\cite{romp}, SWORD~\cite{sword}, and
OMPRacer~\cite{ompracer} specialize happens-before detection for OpenMP. These
are far more powerful race detectors than output differencing, and we use
Archer as an independent instrument rather than as a competitor. Our finding is
not a better detector: it is that for frontier-model output on idiomatic tasks
the race they hunt is largely absent, and the binding constraint has moved to
whether the accepted code is parallel at all.

\subsection{Differential and metamorphic testing}

Executing a program under multiple configurations and flagging disagreements is
the essence of differential testing~\cite{difftesting}; checking properties
preserved under transformation is metamorphic testing~\cite{metamorphic}. The
serial projection is, in this vocabulary, a metamorphic relation run in
reverse: instead of a transform under which the output should be preserved and
a violation indicates a bug, it is a transform under which the output
\emph{is} preserved and that preservation indicates a defect in the reward.

Across all of these lines one gap recurs: measured performance is treated as an
evaluation metric and left out of the objective the model is trained and the
agent is stopped on. Closing that gap - moving performance from the scoreboard
into the reward - is what this paper set out to do, and what we summarize next.
